\documentclass[a4paper, 10pt]{article}

% importazione dei pacchetti
\usepackage[utf8]{inputenc}
\usepackage[T1]{fontenc}
\usepackage[italian]{babel}
\usepackage{xcolor}
\usepackage{graphicx}
\usepackage{tikz}
\usepackage[hidelinks]{hyperref}

% settaggio dei margini
\setlength{\oddsidemargin}{0pt}
\setlength{\evensidemargin}{0pt}
\setlength{\topmargin}{0pt}
\setlength{\headheight}{0pt}
\setlength{\textwidth}{16cm}
\setlength{\textheight}{24cm}

% definizione del comando \cvphoto
\newcommand{\cvphoto}[4]{

    \begin{tikzpicture}[baseline=(pic.base)]

        \clip (0,0) circle (#3/2);
        \node[inner sep = 0pt] (pic) {
                
            \raisebox{-0.45\height}{\includegraphics[width = #3]{#4}}
                
        };
        \draw[line width = #1, draw = #2] (0,0) circle (#3/2);

    \end{tikzpicture}

}

% definizione del comando \cvsection
\newcommand{\cvsection}[1]{

    \vspace{0.5cm} \par \noindent \rule{\textwidth}{0.5pt} \vspace{-0.75cm}
    \section*{#1}

}

% definizione dell'enviroment \cvsubsection
\newsavebox{\cvsubbox}
\newenvironment{cvsubsection}[2]{

    \vspace{0.5cm} \par \noindent
    \begin{minipage}{\textwidth}

        \colorbox{#1}{

            \parbox{\textwidth}{\textbf{#2}}
                        
        }
        \begin{lrbox}{\cvsubbox}

            \begin{minipage}{\dimexpr \textwidth - 0.5em \relax}

}{
    
            \end{minipage}

        \end{lrbox} \par \noindent
        \begin{tabular}{ @{} p{0.5em} | p{\dimexpr \textwidth - 0.5em \relax} @{} }
            
            &   \usebox{\cvsubbox} \\
            
        \end{tabular}
        
    \end{minipage} \par

}

% definizione del comando \cvitem 
\newlength{\cvlabelwidth}
\newcommand{\cvitem}[2]{

    \par \noindent
    \settowidth{\cvlabelwidth}{\textbf{#1:} }
    \hangindent = \cvlabelwidth
    \hangafter = 1
    \textbf{#1:} #2 \par \vspace{0.5em}
    
}

% definizione del comando \skill
\newcommand{\skill}[2]{

  \textbf{#1} \hfill
  \foreach \i in {1,...,5}{
    \begin{tikzpicture}[baseline=-0.5ex]
      \ifnum \i > #2
        \draw[gray!40, fill=gray!10] (0,0) circle (0.3ex);
      \else
        \draw[gray!70, fill=gray!70] (0,0) circle (0.3ex);
      \fi
    \end{tikzpicture}
  }

}

% inizio del documento
\begin{document}

    % Intestazione
    \hspace{-1 cm}
    \begin{tabular}{ @{} c c @{} }
        
        \cvphoto{5pt}{blue!30}{3.5cm}{CV_Photo.jpg} &
        \begin{minipage}{\dimexpr \textwidth - 3.5cm \relax}
            
            \centering
            {\Large \textbf{Lorenzo Di Napoli}} \\[0.5em]
            Data di nascita: 03/12/2004 \quad | \quad Indirizzo: Piazza Bernini 19, Cormano (MI) \\[0.2em]
            Email: lorenzo\_dinapoli@icloud.com \quad | \quad Tel: \texttt{+39 345 599 2100} \\[0.2em]
            Studente di Ingegneria Informatica | \texttt{Politecnico di Milano} \\[0.2em]
            Linkedin: \href{https://www.linkedin.com/in/lorenzo-di-napoli-38108a340}{https://www.linkedin.com/in/lorenzo-di-napoli-38108a340} \\[0.2em]
            Github: \href{https://github.com/LoreDN}{https://github.com/LoreDN} \\[0.2em]

        \end{minipage} \\
        
    \end{tabular} \vspace{0.25cm}

    % Introduzione Personale
    \cvsection{Introduzione Personale}

        Studente di Ingegneria Informatica al Politecnico di Milano (media: 26.11/30), con competenze in programmazione low-level (\textit{VHDL · Assembly Risc-V · C · C++ · CUDA}) 
        e interesse per sistemi embedded e robotica. Orientato alla comprensione profonda dei sistemi, e dell'influenza che l'architettura hardware esercita sul software applicativo.

    % Percorso di studi
    \cvsection{Percorso di Studi}

        \cvitem{2023 - previsto 2026}{Laurea Triennale in Ingegneria Informatica, Politecnico di Milano (MI) --- media attuale: 26.11/30.}
        \cvitem{Corsi Extracurriculari}{Architettura GPU (CUDA) · Architettura FPGA (VHDL \& HLS)}
        \cvitem{2018 - 2023}{Diploma di Maturità presso Liceo Scientifico opzione Scienze Applicate, Istituto Salesiano Sant'Ambrogio, Milano (MI).}

    % Competenze Tecniche & Soft Skills
    \cvsection{Competenze Tecniche \& Soft Skills}
   
        \vspace{-0.5cm}
        \begin{cvsubsection}{blue!20}{Linguaggi di Programmazione}

            \begin{minipage}[t]{0.2\textwidth}
                \skill{C/C++}{5}
                \skill{Java}{4}
            \end{minipage}
            \hspace{1 cm}
            \begin{minipage}[t]{0.2\textwidth}
                \skill{Cuda}{2}
                \skill{Python}{2}
            \end{minipage}
            \hspace{1 cm}
            \begin{minipage}[t]{0.25\textwidth}
                \skill{Assembly}{3}
                \skill{VHDL}{3}
            \end{minipage}

        \end{cvsubsection}
        
        \begin{cvsubsection}{blue!20}{Strumenti e tecnologie di supporto}

            \begin{minipage}[t]{0.3\textwidth}
                \begin{itemize}
                    \item Linux (Ubuntu, Fedora)
                    \item Makefile · CMake
                    \item Docker/Podman
                \end{itemize}
            \end{minipage}
            \hspace{1 cm}
            \begin{minipage}[t]{0.525\textwidth}
                \begin{itemize}
                    \item Git · GitHub
                    \item Valgrind · Cachegrind · Massif · GDB
                    \item Vivado · Vitis
                \end{itemize}
            \end{minipage}

        \end{cvsubsection}
        
        \begin{cvsubsection}{blue!20}{Soft Skills da Esperienze Lavorative}

            \begin{itemize}
                \item Comunicazione efficace e lavoro di squadra
                \item Gestione dell'emergenza e decision-making sotto pressione
                \item Presa della Responsabilità Civile e Penale delle proprie azioni e di quelle altrui
                \item Supervisione e Gestione della sicurezza dei clienti
                \item Abilitazione al Primo Soccorso (BLSD)
            \end{itemize}

        \end{cvsubsection}
        
        \begin{cvsubsection}{blue!20}{Conoscenze Tecniche a partire da corsi universitari}
    
            Architettura dei Calcolatori (RISC-V ) · Reti Logiche · Sistemi Operativi (Linux) \\[0.2 em]
            Algoritmi e Principi dell'Informatica · Ingegneria del Software \\ [0.2 em]
            Fondamenti di Comunicazione e Internet · Database e Sistemi Informativi \\[0.2 em]
            Elettrotecnica · Fondamenti di Elettronica · Fondamenti di Automatica \\[0.4 em]

            Architettura GPU (CUDA) · Architettura FPGA (VHDL \& HLS) \\[0.2 em]
    
        \end{cvsubsection}

    % Esperienze lavorative
    \cvsection{Esperienze lavorative e Progetti}

        \begin{cvsubsection}{blue!20}{Progetti universitari}

            \vspace{0.25 cm}
            \cvitem{\underline{\href{https://github.com/LoreDN/Movhex}{Movhex}}}{sviluppo in \textbf{C} di un programma che consente di creare e modificare una mappa esagonale 2D, con l'obiettivo di trovare il percorso a 
                minor costo tra due punti di essa. \\[0.2 em]
                Il risultato finale ha raggiunto una complessità temporale di $O(n \log (n))$ e spaziale di $O(n)$; dove $n$ è il numero di esagoni presenti nella mappa.}
            \vspace{0.25 cm}
                \cvitem{Priority-List}{design ed implementazione in \textbf{VHDL} di un modulo HW, il quale svolge la funzione di controllore per una Lista con priorità di Task. \\[0.2 em]
                Il risultato finale è stato implementato mediante una FSM totalmente sintetizzabile su FPGA senza Latch.}
                
        \end{cvsubsection}
            
        \begin{cvsubsection}{blue!20}{Progetti personali}
                
            \cvitem{\underline{\href{https://github.com/LoreDN/Cpp}{LDN Library}}}{libreria personale general purpose, implementata in \textbf{C++20} sfruttando le funzionalità avanzate del linguaggio.}
            \cvitem{\underline{\href{https://github.com/LoreDN/Neural-Networks}{Neural Networks}}}{implementazione from-scratch di differenti reti neurali in Python, incentrate sul clustering di diverse figure geometriche 
                (cerchi, spirali, ecc.).}

        \end{cvsubsection}

        \begin{cvsubsection}{blue!20}{Esperienze lavorative}

            \cvitem{Assistente Bagnanti}{Abilitazione al primo soccorso e all'uso del defibrillatore (BLSD), responsabilità civile e penale delle proprie azioni e di quelle dei bagnanti.}
                \begin{itemize}

                    \renewcommand{\labelitemi}{-}
                    \item \textbf{giugno 2025 - dicembre 2025 [116 ore]:} esperienza lavorativa estiva e invernale come Assistente Bagnanti P, svolta presso il centro sportivo \textit{"M.G.M.Sport Srl"} di Paderno Dugnano (MI).
                    \item \textbf{giugno 2024 - luglio 2024 [171 ore]:} esperienza lavorativa estiva come Assistente Bagnanti P, svolta presso il centro sportivo \textit{"M.G.M.Sport Srl"} di Paderno Dugnano (MI).
                    \item \textbf{luglio 2022 [36.5 ore]:} esperienza lavorativa estiva come Assistente Bagnanti P, svolta presso il centro sportivo \textit{"M.G.M. Sport Srl"} di Desio (MB). 
            
                \end{itemize}

        \end{cvsubsection}

    % Conoscenze linguistiche
    \cvsection{Conoscenze Linguistiche}

        \cvitem{Italiano}{Madrelingua.}
        \cvitem{Inglese}{Livello C1, Grado C (Certificazione CAE Cambridge).}
    
    % Certificazioni e Documenti
    \cvsection{Certificazioni e Documenti {\normalsize ---> \underline{\href{https://github.com/LoreDN/LoreDN/raw/main/CV/Certifications.zip}{DOWNLOAD}}}}

        \cvitem{Patente di guida}{Patente di guida di tipo B --- [18/11/2023]}
        \cvitem{Certificazione CAE Cambridge}{Certificazione di lingua inglese, livello C1 Grado C --- [09/01/2026]}
        \cvitem{Brevetto di Assistente Bagnanti}{conseguimento del brevetto di Assistente Bagnanti (livelli P, IP, MIP) presso \textit{"Federazione Italiana Nuoto" (FIN)} --- [09/11/2022]}
        \cvitem{Attestato SFA}{attestato di partecipazione e superamento della \textit{"Scuola Formazione Animatori MGS" (SFA)} presso l'associazione \textit{"Movimento Giovanile Salesiano" (MGS)} --- [17/04/2023]}

\end{document}
